\chapter{Sound \tagZusatz}
\renewcommand{\chaptertitle}{Sound}

\lehead[]{\normalfont\sffamily\hspace*{-2.00cm}\textcolor{white}{\colorbox{lightblue}{\makebox[1.60cm][r]{\thechapter}}}\hspace{0.17cm}\textcolor{lightblue}{\chaptertitle}}
\rohead[]{\textcolor{lightblue}{\chaptertitle}\normalfont\sffamily\hspace*{0.17cm}\textcolor{white}{\colorbox{lightblue}{\makebox[1.60cm][l]{\thechapter}}}\hspace{-2.00cm}}
%\chead[]{}
\rehead[]{\textcolor{lightblue}{AvHG, Inf, My}}
\lohead[]{\textcolor{lightblue}{AvHG, Inf, My}}

\lstset{style=myJava}

Analog zum Laden von Bildern können in Java auch Sound-Dateien geladen und
abgespielt werden.


\section{Wiedergabe von *.wav-Dateien}

Java bringt für unkomprimierte Audiodateien im WAV-Format (16-Bit PCM) bereits
alle nötigen Klassen im Standard-Package \myPackage{javax.sound.sampled} mit:

\begin{lstlisting}
import javax.sound.sampled.AudioSystem;
import javax.sound.sampled.Clip;
import javax.sound.sampled.LineUnavailableException;
import javax.sound.sampled.UnsupportedAudioFileException;

Clip clip = AudioSystem.getClip();
clip.open(AudioSystem.getAudioInputStream(getClass().getResource("ok.wav")));
\end{lstlisting}

Anschließend kann man das \myClass{Clip}-Objekt wahlweise mit
\lstinline|loop(Clip.LOOP_CONTINUOUSLY)| sich ständig wiederholend oder mit
\lstinline|start()| einmalig abspielen lassen.

Mit \lstinline|stop()| lässt sich die Wiedergabe jederzeit anhalten. Der
Programmablauf wartet nicht auf das Ende von \lstinline|start()| bzw.\
\lstinline|loop()|, sondern wird sofort fortgesetzt.

\medskip
\noindent \textbf{Wichtiger Praxistipp:} Wenn ein Soundeffekt (z.\,B.\ ein
Sprung- oder Treffergeräusch) mehrfach hintereinander abgespielt werden soll,
muss der Abspielzeiger vor jedem erneuten \lstinline|start()| wieder an den
Anfang zurückgesetzt werden:

\begin{lstlisting}
clip.setFramePosition(0);   æ// Zurückspulen an den Anfang
æclip.start();
\end{lstlisting}


\subsection{Komplettes Beispiel}

\begin{lstlisting}
import java.awt.EventQueue;
import java.io.IOException;
import javax.sound.sampled.AudioSystem;
import javax.sound.sampled.Clip;
import javax.sound.sampled.LineUnavailableException;
import javax.sound.sampled.UnsupportedAudioFileException;
import javax.swing.JFrame;

public class PlayWAV extends JFrame {

  public PlayWAV() {
    super("PlayWAV");
    setDefaultCloseOperation(JFrame.EXIT_ON_CLOSE);
    setBounds(100, 100, 450, 300);

    æClip sound;
    try {
      sound = AudioSystem.getClip();
      sound.open(AudioSystem.getAudioInputStream(getClass().getResource("ok.wav")));
      sound.start();
      // sound.loop(Clip.LOOP_CONTINUOUSLY);
    } catch (LineUnavailableException | IOException | UnsupportedAudioFileException e) {
      e.printStackTrace();
    }æ
  }

  public static void main(String[] args) {
    EventQueue.invokeLater(() -> {
      try {
        PlayWAV frame = new PlayWAV();
        frame.setVisible(true);
      } catch (Exception e) {
        e.printStackTrace();
      }
    });
  }
}
\end{lstlisting}


\section{Was ist mit *.mp3-Dateien?}

Die Standard-Klassen von Java unterstützen von Haus aus kein MP3-Format (da
dieses lange Zeit patentrechtlich geschützt war). Für die direkte Wiedergabe
von \myFile{*.mp3}-Dateien müsste man externe Zusatzbibliotheken (wie etwa das
JavaFX-Framework, \url{https://openjfx.io} oder Bibliotheken wie \emph{JLayer}) 
aufwendig in das Projekt einbinden.

\medskip
\noindent \textbf{Einfachste Lösung für eigene Projekte:} Wandle vorhandene
MP3-Dateien vorab mit einem Audioprogramm (z.\,B.\ \emph{Audacity} oder dem
\emph{VLC Media Player}) in das Standard-Format \myFile{*.wav} (16-Bit PCM) um.
Gerade für kurze Soundeffekte in Spielen hat \myFile{*.wav} ohnehin den Vorteil,
dass der Ton ohne Verzögerung (Latenz) sofort abgespielt wird.


\section{Bequeme Alternative: Die Klasse HSound}

Da das Laden, Fehlerabfangen und Zurückspulen eines \myClass{Clip}-Objektes
relativ viel Schreibarbeit erfordert, stellt das Package \myPackage{hilfe} die
Klasse \myClass{HSound} bereit. Damit lässt sich eine WAV-Datei mit einer
einzigen Codezeile abspielen -- ganz ohne eigenen \lstinline|try-catch|-Block:

\begin{lstlisting}
import hilfe.*;

// Spielt "ok.wav" aus dem eigenen Package einmalig ab:
HSound.play(this, "ok.wav");

// Spielt "musik.wav" als Endlosschleife ab:
HSound.loop(this, "musik.wav");

// Stoppt die Wiedergabe wieder:
HSound.stop(this, "musik.wav");
\end{lstlisting}

\noindent \textbf{Hinweis:} Der erste Parameter \lstinline|this| teilt der
Klasse \myClass{HSound} mit, in welchem Package nach der Sound-Datei gesucht
werden soll. Falls du dich in einer \lstinline|static|-Methode befindest (in
der es kein \lstinline|this| gibt), kannst du stattdessen den Pfad ab dem
\myFile{src}-Ordner angeben: \lstinline|HSound.play("meinPackage/ok.wav");|.
